\section{Equal-DOF baseline comparisons}
\label{sec:results}

\subsection{Comparison setup and scope}

Five cases use a common evaluator: the one-dimensional linear problem, the one-dimensional cubic
problem, the two-dimensional sinh problem, the one-dimensional pure \(p=4\) problem, and the
two-dimensional pure \(p=4\) problem. The comparison variable is the number of independent global
coefficients: retained atoms for CGA, ridge features held fixed for RFM, and independent
finite-element coefficients after the constant mode is removed where required. Equal degrees of
freedom (DOF) do not imply equal runtime, memory, score evaluations, or nonlinear solver work, and
no cross-method timing conclusion is drawn from runs performed in different environments.
The supplementary numerical data at
\href{https://github.com/RonzolYu/cga4PDE}{\texttt{github.com/RonzolYu/cga4PDE}} give the prescribed realization count, the
available converged realizations, and the solver status for each width.

FEM uses continuous Lagrange elements on uniform one-dimensional meshes or structured triangular
two-dimensional meshes.  P1 and P3 provide low- and high-order mesh-based references in the main
figures. P2 results are retained in the supplementary material. P3 is the most accurate among the tested uniform-mesh FEM variants
at each largest common coefficient count reported below. The linear and cubic FEM curves are close
at high resolution, which is consistent with the common manufactured profile and an \(H^1\) error
dominated by the finite-element approximation.

For each case, the RFM uses a nested family of normalized \(\relu^3\) ridges and fits the outer
coefficients. Ten random-feature realizations are prescribed for each problem and width. Curves show the median and
interquartile range over solver-successful realizations at each width; the prescribed denominator and solver status
are listed in the supplementary numerical records. All ten realizations reach the displayed ranges for the linear,
cubic, sinh, and one-dimensional pure \(p=4\) problems. In the two-dimensional pure \(p=4\) problem,
seven of ten prescribed realizations succeed at \(N=128\) and four of ten at \(N=256\); consequently its
statistical curve ends at 256 and its terminal band uses those successful realizations. No distributional statement
is made at \(N=512\).

The CGA and RFM comparison curves join the computed powers of two listed in the numerical setup and the
terminal retained state needed to define the largest common comparison width, while FEM uses its own mesh degrees of
freedom. The figures contain no common-grid interpolants, theoretical reference lines, or
convergence-order-versus-DOF curves. The reported panel range ends at the terminal CGA width fixed
by the stopping rule, so that endpoint is determined by the stopping rule and the available method
states rather than by the plotted error values. Each method is shown only over the states it
actually reached. The equal-DOF tables use linear interpolation of log(error) against log(DOF)
between adjacent computed values when an exact coefficient count is unavailable; no extrapolation
is used.

The training, evaluation, and refinement quadratures are fixed as specified in the
supplementary protocol; all comparisons use the common DOF set
$8,16,32,64,128,256$, with $512$ used only when every method is bracketed by computed states.

\subsection{Evaluation and statistical samples}

The comparison errors use the fixed evaluation quadrature specified in the supplement, separately from
the quadrature used to construct the features and fit their coefficients. RFM medians and
quartiles are computed for each metric from successful coefficient solves with finite values
of that metric. All ten prescribed realizations contribute at the displayed C1--C4 widths. For C5,
seven of ten solves succeed at $N=128$ and four of ten at $N=256$; unsuccessful solves
remain in the prescribed denominator.

A separate supplementary evaluation contains 617 states satisfying its acceptance and
coefficient-count checks. Its signed energy/metric filter admits 611 of 617 states, its
\(H^1\)-dual residual filter admits 394 of 617, and its two-dimensional Bregman-integration
check admits 137 of 138. These counts refer to that separate collection; they do not define
the samples underlying the RFM curves or endpoint table in this section.


\subsection{Relative Sobolev errors}

For the linear, cubic, and sinh problems,
\Cref{fig:baseline-sobolev-1d,fig:baseline-sobolev-rest} report the relative \(H^1\) error. For the
two pure \(p=4\) problems, the corresponding panels report the relative
\(W^{1,4}\) error. These are Sobolev errors; the natural \(V\)-distance is reported separately below.

\begin{figure}[!htbp]
\centering
\subfloat[Linear, \(d=1\).]{%
\begin{minipage}[t]{0.49\textwidth}
\includegraphics[width=\linewidth]{C1_relative_h1_error.pdf}
\end{minipage}}\hfill
\subfloat[Cubic, \(d=1\).]{%
\begin{minipage}[t]{0.49\textwidth}
\includegraphics[width=\linewidth]{C2_relative_h1_error.pdf}
\end{minipage}}
\caption{Relative \(H^1\) error versus effective degrees of freedom for the one-dimensional
multifrequency problems. RFM curves show medians and interquartile bands over solver-successful realizations
from ten prespecified seeds. All markers are computed states.}
\label{fig:baseline-sobolev-1d}
\end{figure}

\begin{figure}[!htbp]
\centering
\subfloat[Sinh, \(d=2\): relative \(H^1\) error.]{%
\begin{minipage}[t]{0.32\textwidth}
\includegraphics[width=\linewidth]{C3_relative_h1_error.pdf}
\end{minipage}}\hfill
\subfloat[Pure \(p=4\), \(d=1\): relative \(W^{1,4}\) error.]{%
\begin{minipage}[t]{0.32\textwidth}
\includegraphics[width=\linewidth]{C4_relative_w1p_error.pdf}
\end{minipage}}\hfill
\subfloat[Pure \(p=4\), \(d=2\): relative \(W^{1,4}\) error.]{%
\begin{minipage}[t]{0.32\textwidth}
\includegraphics[width=\linewidth]{C5_relative_w1p_error.pdf}
\end{minipage}}
\caption{Relative Sobolev errors for the two-dimensional sinh problem and the pure \(p=4\) problems.
The two-dimensional pure \(p=4\) RFM band uses the seven solver-successful realizations at \(N=128\) and the
four solver-successful realizations at \(N=256\), with ten prescribed seeds at each width.}
\label{fig:baseline-sobolev-rest}
\end{figure}
\FloatBarrier

At the largest common endpoint supported by the configured RFM widths and bracketing FEM states,
CGA has the smallest relative Sobolev error in the linear, cubic, sinh, and one-dimensional pure
\(p=4\) problems; all ten prescribed RFM solves succeed in these four cases. The RFM-median/CGA ratios
are 1.49, 1.52, 1.77, and 5.76, and the P3/CGA ratios are 4.24, 4.35, 5.64, and 1.88. In the
two-dimensional pure \(p=4\) problem the successful-run RFM median and the P3 value are also above the
CGA value, by factors 5.01 and 16.12, respectively. The two-dimensional RFM statistics use seven of ten prescribed realizations at \(N=128\) and four of ten at \(N=256\) after solver-success filtering. These are endpoint comparisons
at equal coefficient count. The curves cross at smaller widths in some cases, so the data do not support
uniform dominance at every DOF.

\begin{table}[!htbp]
\centering\scriptsize
\caption{Relative $H^1$ error for C1--C3 and relative gradient $L^4$ seminorm for C4--C5 at the prescribed common coefficient counts. RFM statistics are median [Q1,Q3] over metric-valid realizations from ten prescribed seeds. Valid means a successful solve, a finite nonnegative metric, and an at-most-one-percent successive-quadrature difference for that metric; incomplete groups are conditional comparisons. Ratios larger than one favor CGA.}
\label{tab:baseline-terminal}
\resizebox{\linewidth}{!}{%
\begin{tabular}{@{}crrrrrr@{}}
\toprule
Problem & DOF & CGA & RFM median [Q1,Q3] & Valid & RFM/CGA & FEM P3/CGA \\ \midrule
Linear, d=1 & 256 & $1.247\!\times\!10^{-4}$ & $1.853\!\times\!10^{-4}$ [$1.779\!\times\!10^{-4}$, $1.913\!\times\!10^{-4}$] & 10/10 & 1.49 & 4.24 \\
Cubic, d=1 & 256 & $1.215\!\times\!10^{-4}$ & $1.853\!\times\!10^{-4}$ [$1.779\!\times\!10^{-4}$, $1.913\!\times\!10^{-4}$] & 10/10 & 1.52 & 4.35 \\
Sinh, d=2 & 256 & $4.074\!\times\!10^{-3}$ & $7.196\!\times\!10^{-3}$ [$6.472\!\times\!10^{-3}$, $8.663\!\times\!10^{-3}$] & 10/10 & 1.77 & 5.63 \\
Pure p=4, d=1 & 128 & $3.395\!\times\!10^{-5}$ & $7.818\!\times\!10^{-5}$ [$6.046\!\times\!10^{-5}$, $1.400\!\times\!10^{-4}$] & 10/10 & 2.30 & 0.88 \\
Pure p=4, d=2 & 256 & $2.493\!\times\!10^{-3}$ & $1.147\!\times\!10^{-2}$ [$9.730\!\times\!10^{-3}$, $1.360\!\times\!10^{-2}$] & 6/10 & 4.60 & 14.54 \\
\bottomrule
\end{tabular}}
\end{table}

\FloatBarrier

\subsection{Natural \texorpdfstring{\(V\)}{V}-distance for the \texorpdfstring{\(p\)}{p}-models}

\Cref{fig:baseline-v} repeats the two pure \(p=4\) comparisons with the relative \(V\)-distance defined in
\Cref{sec:numerics}.  Keeping this geometry separate from the \(W^{1,4}\) error is important because the
two quantities need not have the same terminal behavior.  The figure uses the same computed-state and
RFM uncertainty conventions as the Sobolev panels.

\begin{figure}[!htbp]
\centering
\subfloat[Pure \(p=4,d=1\).]{%
\begin{minipage}[t]{0.49\textwidth}
\includegraphics[width=\linewidth]{C4_v_distance.pdf}
\end{minipage}}\hfill
\subfloat[Pure \(p=4,d=2\).]{%
\begin{minipage}[t]{0.49\textwidth}
\includegraphics[width=\linewidth]{C5_v_distance.pdf}
\end{minipage}}
\caption{Relative \(V\)-distance versus effective degrees of freedom.  No theoretical reference line
or common-grid interpolation is drawn.}
\label{fig:baseline-v}
\end{figure}
\FloatBarrier

\subsection{Interpretation limits}

The equal-DOF results compare finite-pool CGA with the specified FEM and RFM methods.  They do
not establish a runtime ordering, do not compare CGA with every randomized or greedy dictionary method,
and do not turn endpoint ratios into asymptotic complexity estimates. The two-dimensional pure
\(p=4\) RFM summaries use
four of ten solver-successful realizations at the terminal width and are not a complete ten-seed
statistical comparison. Supplementary energy-gap comparisons are more sensitive to quadrature because they
subtract two separately integrated energies. The supplementary numerical data specify the prescribed and
converged realization counts, interpolation brackets, and computed coefficient counts.
